\section{Anhang: Code-Beispiele}\label{sec:anhang}

\subsection{Signal-Slot-Verbindung (PyQt6)}

Das Signal-Slot-Pattern ermöglicht lose Kopplung zwischen Komponenten. Die \texttt{ActionButtons}-Klasse definiert ein Signal, das emittiert wird, wenn ein Spieler eine Aktion wählt:

\begin{lstlisting}[language=Python, caption={Signal-Definition in ActionButtons}]
from PyQt6.QtWidgets import QWidget
from PyQt6.QtCore import pyqtSignal

class ActionButtons(QWidget):
    action_selected = pyqtSignal(str, int)  # action, bet_size
    
    def emit_action(self, action: str, bet_size: int):
        self.action_selected.emit(action, bet_size)
\end{lstlisting}

Die GUI-Klasse verbindet dieses Signal mit einem Handler:

\begin{lstlisting}[language=Python, caption={Signal-Verbindung in AgentVsHumanGUI}]
class AgentVsHumanGUI(AgentVsHumanLayout):
    def setup_connections(self):
        if hasattr(self, 'action_buttons'):
            self.action_buttons.action_selected.connect(self.handle_action)
\end{lstlisting}

\subsection{Strategie-Laden mit Pickle}

Die \texttt{load\_strategy()}-Methode lädt komprimierte Strategie-Dateien mit \texttt{gzip} und \texttt{pickle}:

\begin{lstlisting}[language=Python, caption={load\_strategy() Methode}]
import gzip
import pickle as pkl
from training.solvers.cfr_solver import CFRSolver

def load_strategy(self, filepath):
    with gzip.open(filepath, 'rb') as f:
        data = pkl.load(f)
    
    if 'average_strategy' in data:
        return data['average_strategy']
    else:
        strategy_sum = data['strategy_sum']
        return CFRSolver.average_from_strategy_sum(strategy_sum)
\end{lstlisting}

\subsection{Thread-Safe Game Logic}

\texttt{threading.Lock} wird mit Kontext-Managern verwendet, um Thread-Safety zu gewährleisten:

\begin{lstlisting}[language=Python, caption={Lock-Verwendung in register\_client()}]
import threading
import time
import html
from typing import Optional

class PokerGameLogic:
    def __init__(self, game, game_id=None):
        self.lock = threading.Lock()
        self.clients: dict[int, float] = {}
        self.client_names: dict[int, str] = {}
        self.max_players = 2
    
    def register_client(self, player_id: int, name: Optional[str] = None) -> bool:
        with self.lock:
            self._prune_stale_clients()
            if player_id in self.clients:
                return False
            if len(self.clients) >= self.max_players:
                return False
            self.clients[player_id] = time.time()
            if isinstance(name, str) and name.strip():
                sanitized_name = html.escape(name.strip())
                self.client_names[player_id] = sanitized_name
            return True
\end{lstlisting}

\begin{lstlisting}[language=Python, caption={Lock-Verwendung in handle\_action()}]
def handle_action(self, player_id, action, bet_size):
    with self.lock:
        if not isinstance(player_id, int) or player_id not in (0, 1):
            return False
        if self.game.done:
            return False
        if self.game.current_player != player_id:
            return False
        
        state = self.game.get_state(self.game.current_player)
        legal_actions = state.get('legal_actions', self.game.get_legal_actions())
        
        if action not in legal_actions:
            return False
        
        self.game.step(action)
        self._drain_chance_nodes()
        return True
\end{lstlisting}

\subsection{Rate Limiting Decorator}

Rate Limiting wird als wiederverwendbarer Decorator implementiert:

\begin{lstlisting}[language=Python, caption={Rate Limiting Decorator}]
from functools import wraps
from flask import Flask, jsonify, request
from gui.server.rate_limiter import RateLimiter

class PokerHTTPServer:
    def __init__(self, game, host='localhost', port=8888):
        self.rate_limiter = RateLimiter()
        self._setup_routes()
    
    def _rate_limit(self, max_requests: int, window_seconds: float):
        """Decorator für Rate Limiting."""
        def decorator(f):
            @wraps(f)
            def wrapper(*args, **kwargs):
                ip = request.remote_addr or 'unknown'
                if not self.rate_limiter.is_allowed(ip, max_requests, window_seconds):
                    return jsonify({
                        'status': 'error',
                        'message': 'Rate limit exceeded'
                    }), 429
                return f(*args, **kwargs)
            return wrapper
        return decorator
    
    def _setup_routes(self):
        @self.app.route('/player_id', methods=['GET'])
        @self._rate_limit(max_requests=5, window_seconds=60)
        def get_player_id():
            # ...
\end{lstlisting}

\subsection{WebSocketWorker (Asyncio + QThread)}

\texttt{WebSocketWorker} integriert \texttt{asyncio} mit der Qt Event Loop:

\begin{lstlisting}[language=Python, caption={WebSocketWorker Klasse}]
import asyncio
import json
from PyQt6.QtCore import QThread, pyqtSignal
from websockets.client import connect
from typing import Optional

class WebSocketWorker(QThread):
    state_received = pyqtSignal(dict)
    error_occurred = pyqtSignal(str)
    connected_signal = pyqtSignal(int)
    
    def __init__(self, server_url: str, player_name: Optional[str] = None):
        super().__init__()
        self.server_url = server_url
        self.player_name = player_name
        self.running = False
    
    def run(self):
        """Startet die asyncio-Event-Loop."""
        self.loop = asyncio.new_event_loop()
        asyncio.set_event_loop(self.loop)
        try:
            self.loop.run_until_complete(self._connect_and_listen())
        except Exception as e:
            self.error_occurred.emit(str(e))
        finally:
            self.loop.close()
    
    async def _connect_and_listen(self):
        """Verbindet zum Server und lauscht auf Nachrichten."""
        try:
            self.websocket = await connect(self.server_url)
            self.running = True
            
            # Empfange player_id
            msg = await self.websocket.recv()
            data = json.loads(msg)
            if data.get("type") == "player_id":
                self.player_id = data.get("player_id")
                self.connected_signal.emit(self.player_id)
            
            # Sende Name falls vorhanden
            if self.player_name:
                await self.websocket.send(json.dumps({
                    "type": "register",
                    "name": self.player_name
                }))
            
            # Lausche auf eingehende Nachrichten
            async for message in self.websocket:
                data = json.loads(message)
                if data.get("type") == "state":
                    self.state_received.emit(data.get("state"))
        except Exception as e:
            self.error_occurred.emit(str(e))
\end{lstlisting}

\subsection{pytest-Fixture für Server-Setup}

Eine pytest-Fixture startet den HTTP-Server im Hintergrund:

\begin{lstlisting}[language=Python, caption={pytest-Fixture für HTTP-Server}]
import pytest
import threading
import time
import socket
from gui.server.http_server import PokerHTTPServer
from envs.kuhn_poker.game import KuhnPokerGame

def _free_port():
    with socket.socket(socket.AF_INET, socket.SOCK_STREAM) as s:
        s.bind(("", 0))
        s.setsockopt(socket.SOL_SOCKET, socket.SO_REUSEADDR, 1)
        return s.getsockname()[1]

@pytest.fixture(scope="module")
def http_server_url():
    """Startet HTTP-Server auf freiem Port, liefert Base-URL."""
    port = _free_port()
    game = KuhnPokerGame()
    server = PokerHTTPServer(game, host="localhost", port=port)
    thread = threading.Thread(target=server.start, daemon=True)
    thread.start()
    time.sleep(0.5)
    yield f"http://localhost:{port}"
\end{lstlisting}

\subsection{WebSocket-Broadcast}

Die \texttt{\_broadcast\_state()}-Methode sendet State-Updates an alle verbundenen Clients:

\begin{lstlisting}[language=Python, caption={WebSocket-Broadcast}]
async def _broadcast_state(self):
    """Sendet State an alle verbundenen Clients."""
    for pid, ws in list(self.connections.items()):
        try:
            state = self.game_logic.get_state_update(pid)
            await ws.send(json.dumps({
                "type": "state",
                "state": state
            }))
        except Exception as e:
            logger.error(f"Error sending state to player {pid}: {e}")
            self.connections.pop(pid, None)
            self.game_logic.unregister_client(pid)
\end{lstlisting}
